%======================================================================
% Noise, not optimality: a capacity-to-budget diagnostic for stealthy
% poisoning in differentially private federated learning
% Target: Computer Networks (Elsevier), elsarticle [final,1p,times]
%======================================================================
\documentclass[final,1p,times]{elsarticle}
\journal{Computer Networks}
\usepackage{amsmath,amssymb,amsthm,mathtools}
\usepackage{booktabs,multirow}
\usepackage{graphicx}
\usepackage[colorlinks=true,linkcolor=blue,citecolor=blue]{hyperref}
\usepackage{placeins}
% Relax float placement so figures/tables are not deferred to the end
% of the document (where they would land inside the references).
\renewcommand{\topfraction}{0.85}
\renewcommand{\textfraction}{0.1}
\renewcommand{\floatpagefraction}{0.85}
\setcounter{topnumber}{4}
\setcounter{totalnumber}{6}

\newtheorem{theorem}{Theorem}
\newtheorem{lemma}{Lemma}
\newtheorem{proposition}{Proposition}
\newtheorem{definition}{Definition}
\newtheorem{remark}{Remark}
\newcommand{\etal}{\emph{et al.}}
\newcommand{\opt}{\mathrm{opt}}

\begin{document}
\begin{frontmatter}
\title{Noise, not optimality: a capacity-to-budget diagnostic for
stealthy poisoning in differentially private federated learning}

\author[aqnu_cs]{Xiaolin Zhang}
\author[aqnu_sp]{Lei Lu\corref{cor1}}
\ead{lul@aqnu.edu.cn}
\cortext[cor1]{Corresponding author.}
\affiliation[aqnu_cs]{organization={School of Artificial Intelligence and Computer Science, Anqing Normal University},
  city={Anqing}, postcode={246133}, state={Anhui}, country={China}}
\affiliation[aqnu_sp]{organization={School of Physical Education and Health, Anqing Normal University},
  city={Anqing}, postcode={246133}, state={Anhui}, country={China}}

\begin{abstract}
With client-level differential privacy (DP) at sub-DP noise scales
stacked on coordinate-median
robust aggregation, the deployability of a stealthy poisoning attack is
decided by a single capacity-to-budget ratio $C/\Gamma$: at $K{=}6$,
$\kappa{=}2$ we verify a
deployable phase ($C/\Gamma{=}8.8$) on FashionMNIST and a capacity-bound
phase ($C/\Gamma{=}0.48$) on CIFAR-10. The two-regime stealthy-capacity
law $\eta_{\max}\propto\sigma\, d^{1/4}$, refined on real FL pipelines
(FashionMNIST, CIFAR-10, MNIST), organizes this phase boundary; a
water-filling allocation and a Stackelberg analysis then explain
\emph{why} the transition occurs. Beyond a measurable crossover, stronger
privacy \emph{enlarges} the stealthy attack region and the defender's
equilibrium noise is non-monotone with an interior optimum
($\sigma^\star{=}1.9$ under the central calibration; $\sigma^\star\in[0.4,6.5]$
across plausible privacy valuations), challenging ``more noise is always
safer'' within the calibrated Stackelberg model. Under a backdoor probe
on CIFAR-10 the allocator ordering resolves and the unconstrained
allocators Pareto-dominate the water-filling optimum (on FashionMNIST
every allocator sits on the noise floor), showing allocator optimality
is damage-model dependent; under the median aggregator a closed-form
order-statistic allocation attains the same median shift in the
single-round scalar setting for $4.4$--$1600\times$ less budget, yet its
end-to-end backdoor gain is statistically null. The contribution is
therefore the diagnostic and its mechanism, not allocator dominance:
$C/\Gamma$ is a pre-deployment tool for tuning privacy under poisoning
threats; under central DP applied only to the aggregate the masking
channel vanishes.
\end{abstract}

%======================================================================
\begin{highlights}
%======================================================================
\item The capacity-to-budget ratio $C/\Gamma$ flags a deployable stealth-aware allocation.
\item Two-regime capacity law holds across FashionMNIST, MNIST, and CIFAR-10.
\item Beyond a crossover, stronger privacy enlarges the stealthy attack surface.
\item Defender utility in noise scale is non-monotone with a calibrated interior optimum.
\item No allocator dominates; the contribution is the $C/\Gamma$ diagnostic, not allocator dominance.
\end{highlights}

\begin{keyword}
federated learning, differential privacy, poisoning attacks, robust
aggregation, capacity-to-budget diagnostic, Stackelberg games
\end{keyword}

\begin{graphicalabstract}
\includegraphics[width=\linewidth]{figures/graphical_abstract/graphical_abstract.pdf}
The joint defense stacks client-level DP noise on coordinate-median robust
aggregation. The capacity-to-budget ratio $C/\Gamma$ (center) converts this
interaction into a single pre-deployment diagnostic: at $K{=}6$,
$\kappa{=}2$ it is deployable on
FashionMNIST ($C/\Gamma{=}8.8$) yet capacity-bound on CIFAR-10
($C/\Gamma{=}0.48$). The defender (left) and the budget-bounded attacker
(right) are coupled through $C/\Gamma$, which organizes every experiment.
\end{graphicalabstract}

\end{frontmatter}

%======================================================================
\section{Introduction}
%======================================================================
Defenders who deploy federated learning (FL) increasingly stack two
protections at once: differential privacy (DP) noise for privacy and
robust aggregation for poisoning tolerance. A natural pre-deployment
question is therefore: \emph{does adding privacy help or hurt the
defender against a stealthy poisoner?} We answer it with a single
deployability diagnostic, the capacity-to-budget ratio $C/\Gamma$, and
show the answer is dataset-dependent: under client-level DP with
coordinate-median aggregation the attack is deployable on FashionMNIST
($C/\Gamma{=}8.8$) yet capacity-bound on CIFAR-10 ($C/\Gamma{=}0.48$;
both at $K{=}6$, $\kappa{=}2$).
This journal has repeatedly studied the same privacy--robustness
coupling from the \emph{defender's} side
\cite{zhou2024roppfl,liu2025antidotefl,zhao2025fedmp}; our contribution
is to approach it from the attacker's resource-allocation side and turn
the interaction into a quantified, pre-deployment phase boundary.

Federated learning \cite{mcmahan2017,kairouz2021advances,li2020federated}
trains shared models across holders whose raw data cannot be pooled. Each
protection is justified on its own; together they interact in a way neither
was built for: DP noise
\cite{dwork2006calibrating,abadi2016,mcmahan2018dp,geyer2017dp} blurs
precisely the signal that robust aggregation
\cite{blanchard2017,yin2018,cao2021fltrust,nguyen2022flame} relies on to
flag a poisoned update, so a defender who raises the noise scale to
``improve privacy'' can inadvertently widen the attack surface---an
interaction that recent studies confirm qualitatively
\cite{mothukuri2021survey,hossain2025mpelm,wang2025ldpfl}. For the defender, this coupling is
first of all a \emph{deployment decision}: before releasing a pipeline
she must fix the noise scale $\sigma$ and the defense budget $\kappa$,
and she needs to know before training whether that specific
configuration still leaves room for a stealthy poisoner. We show
that this question reduces to a single pre-deployment check: the
capacity-to-budget ratio $C/\Gamma$, computed from quantities the server
can already measure, separates a deployable phase from a capacity-bound
one. $C/\Gamma$ is therefore the defender's own go/no-go gauge for
tuning privacy under poisoning threats, not a by-product of attack-side
analysis. We make it quantitative: the
analysis below shows how far the detection separation of a poisoned
update collapses as the noise scale $\sigma$ grows, and what a
budget-bounded attacker should optimally do about it (Fig.~\ref{fig:1}).

\emph{Scope.} All results concern \emph{per-update (client-level)}
DP-style Gaussian noise at sub-DP scales, under a strongest-information
attacker at a malicious or compromised server (client-side adversaries
are upper-bounded by it); \S\ref{sec:setting} gives the formal scope
statement and its central-DP boundary.

Optimal poisoning as constrained optimization has an established
lineage outside FL \cite{biggio2012poisoning,meizhu2015teaching}. On
the attacker side of FL, existing results are either heuristic (noise
used as cover, magnitude chosen by hand) or DP-unaware (optimal
attacks against robust aggregation only)
\cite{shejwalkar2021turning,bhagoji2019analyzing,fang2020local}. None
answers the attacker-side counterpart: \emph{given a bounded attack
budget, what is the optimal allocation under the joint defense?} This
paper does.

\subsection{Contributions}
\begin{enumerate}
  \item[\textbf{C1}] \textbf{A deployability diagnostic for the
  privacy--robustness coupling.} We introduce the capacity-to-budget
  ratio $C/\Gamma$ and show, on real FL pipelines, that it decides
  whether a stealth-aware allocation can deploy at all: $C/\Gamma{=}8.8$
  (deployable) on FashionMNIST versus $C/\Gamma{=}0.48$ (capacity-bound)
  on CIFAR-10. $C/\Gamma$ organizes every experiment below as a single
  phase diagram---deployable when $C\ge\Gamma$, capacity-bound when
  $C<\Gamma$---giving defenders a pre-deployment check rather than a
  per-attack analysis.
  \item[\textbf{C2}] \textbf{A two-regime stealthy-capacity law (T1).}
  The maximum per-client stealthy perturbation obeys
  $\eta_{\max}(\sigma)=\Theta(\sigma\, d^{1/4})$ in the
  isotropic-statistics regime against coordinate-median aggregation, a
  correction to the intuitive $\sqrt{d/d_0}$ heuristic, verified by
  data collapse ($\mathrm{sep}=0.30\,\rho^2/(\sigma^2\sqrt{d})$,
  $R^2=0.984$, Fig.~\ref{fig:4}). On a real FL pipeline with structured benign
  heterogeneity we further derive a \emph{two-regime refinement} in which
  $\eta_{\max}$ saturates into a heterogeneity floor and grows only
  sublinearly in $\sigma$ (Fig.~\ref{fig:7}). The per-client capacities
  $\eta_{\max,i}$ supplied here define the aggregate capacity $C$.
  \item[\textbf{C3}] \textbf{Why the transition happens (mechanism).}
  The stealth-constrained allocation is water-filling (T2) and the
  defender's equilibrium noise is non-monotone with an interior optimum
  $\sigma^\star(a)\in[0.4,6.5]$ across plausible privacy valuations
  ($\sigma^\star{=}1.9$ at the central calibration; T3--T4), challenging
  ``more noise is always
  safer'' within the calibrated Stackelberg model; $72\%$
  ($66$--$76\%$ under $\pm20\%$ slope perturbation) of the marginal over-privatization cost is borne by
  \emph{attack-surface growth}, not the defender's own accuracy
  (Fig.~\ref{fig:2}, Fig.~\ref{fig:5}). Water-filling and the Stackelberg
  analysis are the tools that \emph{explain} the $C/\Gamma$ phase
  boundary, not the contribution itself; their optimality is
  damage-model-relative by construction---the objective is a linear
  influence proxy, whereas realized damage under the deployed median
  aggregator is governed by order statistics (\S\ref{sec:median-alloc}).
\end{enumerate}

%======================================================================
\section{Related Work}
%======================================================================
\emph{Optimal attacks against robust aggregation (no DP).} Poisoning
as constrained optimization predates FL
\cite{biggio2012poisoning,meizhu2015teaching,biggio2018wild}; within
FL, adversarial-optimization \cite{bhagoji2019analyzing},
direction/magnitude optimization against robust aggregators
\cite{shejwalkar2021turning}, and Krum/median-targeted attacks
\cite{fang2020local,sun2019can}, plus backdoor variants
\cite{bagdasaryan2020backdoor,xie2020dba}---all DP-free, with no
masking-capacity model and no per-client allocation question.
\emph{Robust aggregation and detection.} Krum \cite{blanchard2017},
median/trimmed mean \cite{yin2018}, Bulyan \cite{mhamdi2018hidden},
signSGD \cite{bernstein2018signsgd}, trust bootstrapping
\cite{cao2021fltrust}, clustering filters \cite{nguyen2022flame}, and
filter variants \cite{xu2024minvar,xu2024agnostic,luong2024fedlsae}
define the detection layer our attacker must defeat; secure aggregation
\cite{bonawitz2017practical} is orthogonal (our attacks are
server-visible by construction). \emph{Attacks exploiting DP noise
(heuristic).} DP hardening can amplify vulnerability
\cite{song2019privacy}; Hossain \etal\ \cite{hossain2025mpelm} use DP
noise as disguise and Wang \etal\ \cite{wang2025ldpfl} craft
LDP-adaptive attacks---both empirical, and neither addresses budget
allocation; the latter targets LDP rather than central DP.
\emph{Budget-constrained poisoning outside FL.} Qiu \etal\
\cite{qiu2025budget} allocate a global budget via offline RL in
(non-FL) training; per-client stealthy capacities and DP modulation
are absent. \textbf{Positioning.} Unlike these efforts, we are, to our
knowledge, the first to pose
attacker-side \emph{resource allocation} as constrained optimization
under the \emph{joint} DP + robust-aggregation defense, and the first
to derive the resulting safety threshold and Stackelberg
equilibrium.

%======================================================================
\section{System and Threat Model}
%======================================================================
\subsection{Setting}\label{sec:setting}
$N$ clients train jointly via federated averaging~\cite{mcmahan2017,kairouz2021advances};
benign update of client $i$ lies near a $d_0$-dimensional
manifold with isotropic floor noise; the server protects each client
update with client-level Gaussian DP~\cite{dwork2014algorithmic,mcmahan2018dp}: each
(clipped) update is perturbed by $\mathcal{N}(0,\sigma^2 I_d)$ before
aggregation (a client-level DP mechanism: per-update noising before
aggregation), and the server aggregates with a robust rule
(coordinate median~\cite{yin2018,huber1964robust}; extensions in \S\ref{sec:aggr-swap}). The server
flags clients whose distance from the \emph{noised} aggregate exceeds a threshold
calibrated to FPR $5\%$; because the noise enters every client update, both the
per-client statistics and the aggregate are perturbed, which is why the
detector's separation falls with $\sigma$ (the masking effect studied here). For the
Gaussian mechanism with L2 sensitivity $\Delta$ (the full per-update
norm) and per-coordinate noise scale $\sigma$, the analytic Gaussian
mechanism gives $\varepsilon = 2\Delta^2\ln(1.25/\delta)/\sigma^2$ at a
fixed $\delta$~\cite{dwork2006calibrating}; in our pipeline the noise is
calibrated to the \emph{per-coordinate} detection scale rather than to a
meaningful per-update DP guarantee, so we demonstrate the detection-masking
\emph{mechanism} at sub-DP noise scales: in our pipeline
$\Delta\approx0.36$ and $\sigma=\sigma_{\mathrm{rel}}\cdot0.0019$, so the
studied $\sigma_{\mathrm{rel}}\in\{0.5,\dots,16\}$ give per-round
$\varepsilon\gg1$---DP-style Gaussian noise, not a meaningful end-to-end
DP guarantee (the $\varepsilon(\sigma)$ values
and the central-DP boundary argument are in the Supplementary Material).
Under \emph{true} central-DP-on-the-aggregate the server could detect on
the clean pre-noise aggregate and the masking term would vanish, but robust
aggregators are nonlinear, so DP-on-the-aggregate does not compose with
them---whereas deployments that stack DP with robust aggregation must
noise client updates before the aggregator, as in client-level DP
training~\cite{mcmahan2018dp}; we flag the boundary explicitly wherever a
conclusion depends on it (see the Conclusion). We measure
detectability by the Fisher
separation $\mathrm{sep}=(\mu_{\mathrm{mal}}-\mu_{\mathrm{ben}})/s_{\mathrm{pooled}}$
with critical level $\mathrm{sep}^{*}=2$. This is the separation at which a
benign-distribution $5\%$-FPR detector still flags the poisoned client with
probability $1-\Phi(2-1.645)\approx 36\%$; we therefore use
$\mathrm{sep}^{*}=2$ as a \emph{design} separation for the $5\%$-FPR detector
and report it as a stealth cap relative to that benign-design threshold, not as
a zero-detection guarantee. The realized per-allocator operating points (ROC /
true-positive rate) are reported in \S\ref{sec:real}.

\subsection{Attacker}\label{sec:attacker}
The attacker controls $K$ clients, chooses the poisoned set $S$ ($|S|\le K$),
directions, and magnitudes $p_i$, subject to
\begin{equation}
  \sum_{i\in S} p_i^2 \;\le\; \Gamma, \qquad
  0\le p_i \;\le\; \eta_{\max,i}(\sigma),
  \label{eq:budget}
\end{equation}
and maximizes hidden impact $\sum_{i\in S}\alpha_i p_i$ over
non-detected injections. Here $\alpha_i$ is client $i$'s \emph{leverage}:
the first-order (mean-type, FedAvg-style) influence weight of client $i$'s
injected norm on the aggregate loss, calibrated per client from the
real-pipeline probe (e.g.\ $\alpha_i\sim\mathrm{LogN}(0,0.8)$; \S\ref{sec:r1group}).
We use $\sum_i\alpha_i p_i$ as a \emph{linear} influence proxy for
tractability: under the coordinate-median aggregator, realized damage is
governed by order statistics (derived in \S\ref{sec:median-alloc}), so
the proxy is a monotone upper envelope of real impact; we treat it as
the allocator's objective and report realized (median-aggregator)
damage separately in \S\ref{sec:real}.
The attacker estimates $\{\eta_{\max,i}\}$ by probing historical
rounds (a one-shot offline read of past aggregated
updates), with the attacker's adaptive capability bounded by this estimate
(cf.\ the attacker-knowledge assumption in Limitations). Concretely,
the read requires the per-round \emph{aggregated} updates
$\Delta_t=w_t-w_{t-1}$---public to every participant---plus per-client
attribution of $\{\eta_{\max,i}\}$, which we locate at a malicious or
compromised server (the natural reader of all client updates) rather
than at a client-side adversary; a client-side attacker with weaker
information is bounded above by this strongest-information attacker,
and the Supplementary Material shows that even the strongest gains nothing from
reactivity.

\begin{figure}[!t]
\centering
\includegraphics[width=\linewidth]{figures/fig1_overview.pdf}
\caption{Reader map of the joint-defense / resource-allocation problem.
\emph{Defender} (left) stacks server-side DP noise $\sigma$ with
coordinate-median robust aggregation and a $5\%$-FPR detector; the same
noise that protects privacy also \emph{masks} poisoned updates from the
detector (detection separation $\mathrm{sep}$ falls as $\sigma$ rises).
\emph{Attacker} (right) controls $K$ clients under a fixed budget
$\Gamma$ and must choose which clients to poison and how much to inject,
subject to the stealth cap $\eta_{\max,i}(\sigma)$. The two sides are
coupled by $\sigma$: both the attacker's stealth room and the defender's
detection power scale with the same noise, which is why ``more privacy''
can help the attacker. This trade-off is the subject of the paper.}
\label{fig:1}
\end{figure}

%======================================================================
%======================================================================
\section{Implementation and Deployment Considerations}
%======================================================================
The diagnostics above are only useful if a defender can compute them from
observable quantities. We address the three questions a deployer asks.

\paragraph{Attacker information cost}
The attacker estimates the per-client stealth ceilings
$\{\eta_{\max,i}\}$ by a one-shot offline read of historical aggregated
updates $\Delta_t=w_t-w_{t-1}$---public to every participant---and
attributes each scalar $\eta_{\max,i}$ to client $i$
(\S\ref{sec:attacker}). Reading the last $T{=}50$ rounds of model
updates costs $O(Td)$ storage (a few MB for the TinyCNN pipeline), plus a
single pass to fit the benign/malicious separation model. We locate this
strongest-information attacker at a malicious or compromised server
(\S\ref{sec:attacker}); a client-side adversary with strictly weaker
information is therefore upper-bounded by it, so every quantitative claim
is conservative with respect to the attacker's knowledge.

\paragraph{Computing $C/\Gamma$ at runtime}
The aggregate capacity $C{=}\sum_{i\in S}\eta_{\max,i}^2$ needs the
attacker's chosen client set $S$, which is a \emph{threat-model}
parameter, not an observable. The exact worst case over all sets of
size $K$ is the sum of the $K$ largest per-client ceilings,
\begin{equation}\label{eq:cub}
C_{\mathrm{ub}}(K)=\textstyle\sum_{i=1}^{K}\eta_{\max,[i]}^2,
\qquad C\le C_{\mathrm{ub}}(K)\ \ \text{for every } |K'|=K,
\end{equation}
a strict bound because no $K$-set sum can exceed it; a budget-constrained
attacker can only do worse. The ceilings themselves are estimated from
the benign pool, so $C_{\mathrm{ub}}$ inherits, but does not amplify,
that estimation error; ranking clients by leverage $\alpha_i$ is a
computable \emph{heuristic} for selecting the top-$K$ set, not an upper
bound in general---it is near-exact when $\eta_{\max,i}$ varies mildly across clients, as on
FashionMNIST (where $\eta_{\max,i}$ is set primarily by $\sigma$ and
$d$); on CIFAR-10 the per-client floor $h_i$ varies widely (so
$\eta_{\max,i}$ spans a $25\times$ range and is anticorrelated with
$\alpha_i$), and the exact $C_{\mathrm{ub}}$ of Eq.~\ref{eq:cub} should
be used there. On CIFAR-10 even
$C_{\mathrm{ub}}$ leaves $C/\Gamma{<}1$ (Table~\ref{tab:r2cifar}); the
capacity-bound verdict is therefore robust to the defender's ignorance
of $K$ and can be issued before any attack is observed.

\paragraph{Calibrating the detection threshold without attacker labels}
The $5\%$-FPR detector uses the critical separation
$\mathrm{sep}^*{=}2$, i.e.\ a per-client threshold at the benign
$95$th percentile $\mu_{\mathrm{ben}}{+}1.645\,\sigma_{\mathrm{ben}}$
estimated from the benign update pool alone---no attacker labels are
required. Because both the benign and malicious distributions shift with
the DP scale $\sigma$, the threshold must be re-fitted per $\sigma$ (in
practice from a rolling benign window); the separation-based formulation
makes this recalibration a one-line change rather than a retraining step.

\paragraph{Defender's takeaway}
In practice $C/\Gamma$ is used as a sweep, not a single number: fix the
pipeline and sweep the defense budget $\kappa$ over a range that spans
the deployable regime (e.g.\ $\kappa{=}0.25$) and the capacity-bound
regime (e.g.\ $\kappa{=}2$), computing $C_{\mathrm{ub}}(K)$
(Eq.~\ref{eq:cub}) and $\Gamma$ at each point. The phase boundary sits
where $C/\Gamma$ crosses $1$; its location in $\kappa$ tells the
defender how much budget buys a capacity-bound pipeline and how far the
current configuration sits from that boundary. Operate comfortably on
the capacity-bound side; if the boundary is unreachable at an acceptable
$\kappa$, treat the deployable-phase verdict as a signal to shrink $K$
(client screening, attestation) rather than to raise the noise further.

%======================================================================
\section{Stealthy Capacity and Optimal Allocation}
%======================================================================
\subsection{T1: masking-modulated capacity}
\begin{theorem}[Scaling law]\label{thm:t1}
For coordinate-median detection with Gaussian DP noise of scale
$\sigma$ in dimension $d$, the Fisher separation of a norm-$\rho$
isotropic attack satisfies
$\mathrm{sep}(\rho,\sigma,d)=c\,\rho^2/(\sigma^2\sqrt{d})+o(\cdot)$ in the
noise-dominated regime where the benign heterogeneity floor is
negligible ($h\approx0$, so $w^2=\sigma^2$; cf. Lemma~\ref{lem:floor}),
where $c$ is the standardized detection constant
(Remark~\ref{rem:constant}) and $\mathrm{sep}^{*}$ is the design
separation of \S\ref{sec:setting}; hence the largest stealthy magnitude
is
\begin{equation}
  \eta_{\max}(\sigma)\;=\;\sqrt{\tfrac{\mathrm{sep}^{*}}{c}}\;\sigma\, d^{1/4},
  \label{eq:etamax}
\end{equation}
increasing linearly in $\sigma$ and only as $d^{1/4}$ in $d$, and
insensitive to $d_0$ once the benign floor noise covers the orthogonal
subspace.
\end{theorem}
\emph{Proof.} See the Supplementary Material. \hfill$\square$

\begin{lemma}[Floor-noise sufficiency and dimension insensitivity]
\label{lem:floor}
Let benign updates be $x_b=m+\xi$ with per-coordinate second moment
$\mathbb{E}[\xi_j^2]=w^2$ and consider the coordinate-median distance
statistic $s=\|x-m\|^2$ for an attacker displacement $\Delta$,
$\|\Delta\|=\rho$. Then (i)
$\mathbb{E}[s\mid\Delta]=\mathbb{E}[s]+\rho^2$, \emph{independent of
the direction of $\Delta$}; (ii) the benign standard deviation of $s$
is $\sqrt{2}\,w^2\sqrt{d}+o(\cdot)$; hence
\begin{equation}
  \mathrm{sep}(\rho)=c\cdot\frac{\rho^2}{w^2\sqrt{d}},
  \qquad w^2=\sigma^2+h^2,
  \label{eq:sep2}
\end{equation}
where $c$ is the standardized detection constant and $h$ is the
benign-heterogeneity floor per coordinate. In particular, the
intrinsic dimension $d_0$ of the benign manifold enters \eqref{eq:sep2}
only through the estimable floor $h$, not directly (a delta-method
prediction; the direct dimension sweep in Fig.~\ref{fig:4} is isotropic,
so it confirms the $\sqrt{d}$ scaling of the benign spread without
separately isolating intrinsic from ambient dimension).
\end{lemma}
\emph{Proof.} Expanding the square, the cross term has zero mean under
symmetric benign dispersion, so the shift is exactly $\rho^2$
regardless of how $\Delta$ splits across the manifold and its
orthogonal complement; $s/w^2$ is a sum of $d$ standardized squares
(a $\chi^2_d$ scale family) with standard deviation $\sqrt{2d}\,w^2$.
\hfill$\square$

\begin{remark}[Constant alignment]\label{rem:constant}
For the pure Gaussian surrogate with a squared-distance statistic,
$c=1/\sqrt{2}\approx0.707$; the measured $c=0.30$
(Fig.~\ref{fig:4}) decomposes as $c=\gamma/\sqrt{2}$ with
$\gamma\approx0.42$, a median-absorption and finite-sample factor
depending on the client population. The functional \emph{form} (what
Theorem~\ref{thm:t2} consumes via the measured $\eta_{\max,i}$) is
pipeline-independent; the constant is not, so the defender
calibrates $\eta_{\max,i}$ empirically rather than from
\eqref{eq:etamax} alone.
\end{remark}

\emph{Empirical verification (isotropic regime).} Across
$\sigma\in\{2,4,8\}$, $d\in\{600,1500,3000\}$ and six magnitudes each
(54 configurations), all measurements collapse onto
$\mathrm{sep}=0.30\,x$ with $x=\rho^2/(\sigma^2\sqrt{d})$
($R^2=0.984$); Fig.~\ref{fig:4}. Measured $\eta_{\max}/\sigma\in[14.1,18.5]$
and $\eta_{\max}/d^{1/4}\in[9.1,9.4]$ (Fig.~\ref{fig:4}B).

\begin{figure}[!t]
\centering
\includegraphics[width=0.9\linewidth]{figures/fig4_collapse.pdf}
\caption{T1 scaling. \textbf{A}: data collapse: Fisher separation vs.\
$x=\rho^2/(\sigma^2\sqrt{d})$ over $54$ isotropic configurations
($\sigma\in\{2,4,8\}$, $d\in\{600,1500,3000\}$, six magnitudes each);
the through-origin fit $\mathrm{sep}=0.30\,x$ ($R^2=0.984$).
\textbf{B}: $\eta_{\max}$ vs.\ dimension $d$ follows the $d^{1/4}$ law,
refuting the $\sqrt{d/d_0}$ heuristic. Together the panels support the
$\eta_{\max}\propto\sigma\, d^{1/4}$ scaling law.}
\label{fig:4}
\end{figure}

\begin{proposition}[Two-regime capacity law]\label{prop:tworegime}
With a benign-heterogeneity floor $h$ per coordinate that does not
scale with $\sigma$, inverting \eqref{eq:sep2} at the design level
$\mathrm{sep}(\rho)=\mathrm{sep}^{*}$ (the same critical level as in
Theorem~\ref{thm:t1}) yields the stealthy capacity
\begin{equation}
  \eta_{\max}(\sigma)=\sqrt{\tfrac{\mathrm{sep}^{*}}{c}}\;
  d^{1/4}\sqrt{\sigma^{2}+h^{2}}
  \;=\;\eta_{\mathrm{f}}\sqrt{1+(\sigma/\sigma_c)^{2}},
  \label{eq:tworegime}
\end{equation}
with crossover $\sigma_c:=h$ and floor
$\eta_{\mathrm{f}}:=\sqrt{\mathrm{sep}^{*}/c}\;h\,d^{1/4}$; here $c$
and $\mathrm{sep}^{*}$ are the same standardized detection constant and
design separation as in Theorem~\ref{thm:t1}---the floor $h$ enters
only through $w^2=\sigma^2+h^2$ in \eqref{eq:sep2}, not through $c$.
For $\sigma\gg\sigma_c$ the $\sigma\,d^{1/4}$ law of
Theorem~\ref{thm:t1} is recovered (noise-dominated); for
$\sigma\ll\sigma_c$ capacity saturates at the $\sigma$-independent
floor $\eta_{\mathrm{f}}$ (heterogeneity-dominated).
\end{proposition}

\paragraph{Real-pipeline refinement}\label{sec:t1-refine}
The $\propto\sigma\, d^{1/4}$ law describes the
\emph{isotropic-statistics} regime (benign updates drawn like the noise).
On a real FL pipeline (FashionMNIST, a $d{=}20490$ CNN, coordinate-median
server with central DP noise, $60$ warm-up rounds, test accuracy $0.846$),
probing the per-client stealthy capacity yields a rising profile
(Fig.~\ref{fig:7}, left; Table~\ref{tab:r1}): the ratio
$\eta_{\max}/\sigma$ falls from $3.69$ to $0.15$, so the law is not
$\propto\sigma$ in the low-$\sigma$ regime. The cause is a
structured benign-heterogeneity \emph{floor}: the detection spread is set
by benign update variance, which does not grow with $\sigma$, while the
attack's noncentral shift shrinks as the noise inflates
$\mathbb{E}[s]$. Real-pipeline capacity therefore follows the
\emph{two-regime} law \eqref{eq:tworegime}: fitting it to all ten
probes yields $\eta_{\mathrm{f}}=1.96$ and $\sigma_c=13.4$ (in
$\sigma_{\mathrm{rel}}$ units; $R^2=0.997$), with the floor inside the
observed plateau ($\sigma_{\mathrm{rel}}\le 8$) and the crossover
predicting the transition between $\sigma_{\mathrm{rel}}=8$ and $16$.
\paragraph{Unit convention and identification caveat}
The fitted $\sigma_c=13.4$ and the model parameter $h$ in
Proposition~\ref{prop:tworegime} ($\sigma_c:=h$) are the same quantity,
in per-coordinate $\sigma_{\mathrm{rel}}$ units; the apparent gap with
the benign per-coordinate standard deviation ($\approx 1$) is resolved
by the norm conversion carried inside the floor $\eta_{\mathrm{f}}$
(a benign-update-norm quantity). On the rising branch
($\sigma_{\mathrm{rel}}\ge 24$) $\eta_{\max}/\sigma$ plateaus at
$0.15\pm0.005$: the $\propto\sigma\, d^{1/4}$ law of
Theorem~\ref{thm:t1} is \emph{recovered} once $\sigma\gg\sigma_c$, and
$R^2{=}0.997$ over the full window fits the two-regime form
tightly. The defender optimization in \S\ref{sec:stackelberg} uses
this refined form.

\subsection{T2: water-filling allocation}
\begin{theorem}[Stealth-constrained allocation: water-filling form]\label{thm:t2}
For the problem
$\max_{S,p}\sum_{i\in S}\alpha_i p_i$ subject to \eqref{eq:budget}, the
per-client form
\begin{equation}
  p_i^\star=\min\Bigl\{\eta_{\max,i},\ \tfrac{\alpha_i}{2\mu}\Bigr\},
  \qquad \mu\ge 0,\quad
  \sum_{i\in S}(p_i^\star)^2=\min\{\Gamma,\,C\},\quad
  C:=\sum_{i\in S}\eta_{\max,i}^2,
  \label{eq:wf}
\end{equation}
where $\mu\ge 0$ is the Lagrange multiplier on the budget constraint
chosen so that the funded mass matches $\min\{\Gamma,C\}$ (with
$\mu{=}0$ iff $C<\Gamma$), is optimal for a \emph{fixed} funded set $S$. The set $S$ and whether
the cap binds are regime-dependent:
\begin{enumerate}
\item \textbf{Capacity-limited ($C\!<\!\Gamma$).} The budget can never be
  exhausted, so the optimum is the trivial $p_i^\star=\eta_{\max,i}$
  for all funded clients ($\mu{=}0$), with objective
  $\sum_{i\in S}\alpha_i\eta_{\max,i}$.
\item \textbf{Capacity-available ($C\!\ge\!\Gamma$).} The budget binds
  ($\sum(p_i^\star)^2{=}\Gamma$). When every funded client is
  cap-saturated the greedy ordering by decreasing
  $\alpha_i\eta_{\max,i}$ is optimal; when no cap binds the ordering
  reduces to decreasing $\alpha_i$. The unconditional claim ``order by
  $\alpha_i\eta_{\max,i}$ for all $\Gamma$'' is therefore only the
  high-budget limit and is not universally valid.\footnote{It fails
  when caps bind asymmetrically: a low-$\alpha$ client carrying a much
  larger cap can be capped while a higher-$\alpha$ client is not
  (e.g.\ $\alpha{=}(10,1)$, $\eta_{\max}{=}(1,100)$, $\Gamma{=}1$,
  $K{=}1$: selecting the higher-$\alpha$ client dominates).}
\end{enumerate}
\end{theorem}
\emph{Proof.} See the Supplementary Material. \hfill$\square$

The \emph{capping} term is the qualitative difference from
DP-unaware allocation \cite{qiu2025budget}: budget that cannot be
hidden on high-leverage clients is redirected, not lost.

\subsubsection{Median-aggregator realized damage: the order-statistic allocation}
\label{sec:median-alloc}
The objective in \eqref{eq:budget} ($\max\sum_i\alpha_i p_i$) is a
\emph{mean-type} (FedAvg-style) influence proxy. Under the
coordinate-median aggregator actually deployed, the true damage is an
\emph{order-statistic} function of the injected norms, not their sum;
we make the caveat of \S\ref{sec:attacker} quantitative.

\emph{Model.} Fix the attack direction $v$; the benign projections
$u_i=\langle U_i,v\rangle$ (observed via the offline read of
\S\ref{sec:attacker}) are ordered $u_{(1)}\le\cdots\le u_{(N)}$, with
median $m_0=u_{(m)}$, $m{=}\lceil N/2\rceil$. The attacker adds
$p_i\ge0$ on a subset $S$ ($|S|\le K$); the poisoned output is
$g=\mathrm{median}_i(u_i+\delta_i p_i)$, $\delta_i=\mathbf{1}[i\in
S]$, and the realized damage is $D(S,p)=g-m_0$. Crossing a threshold
$t$ requires pushing the $r(t){=}\max\{0,\,m{-}|\{i:u_i\ge t\}|\}$
lowest clients above $t$, each exactly by $t-u_i$ (any excess raises
$\sum p_i^2$ without raising $g$), so the minimum budget to achieve
median $\ge t$ is
\begin{equation}
B(t)=\sum_{i\in S^\star(t)}(t-u_i)^2,\qquad
S^\star(t)=\text{the }r(t)\text{ eligible clients with }u_i<t
\text{ and largest }u_i,
\label{eq:crosscost}
\end{equation}
where \emph{eligible} means $t-u_i\le\eta_{\max,i}$ (otherwise the
injection is detected and the client removed).

\begin{theorem}[Order-statistic allocation]\label{thm:t2median}
The median-shift optimum is
\begin{equation}
p_i^\star=\min\{\eta_{\max,i},\,(t^\star-u_i)_+\},\qquad
t^\star=\max\{t: B(t)\le\Gamma\},
\label{eq:medianalloc}
\end{equation}
with $S^\star$ the eligible pivotal clients just below $t^\star$: only
clients in the order-statistic band $(m_0,\,t^\star]$ receive positive
budget, each exactly the crossing amount $(t^\star-u_i)$. This is
structurally different from the mean-type water-filling of
Theorem~\ref{thm:t2}, which spreads budget in proportion to $\alpha_i$
and never exhibits a threshold-crossing form.
\end{theorem}
\emph{Proof.} See the Supplementary Material. \hfill$\square$

\emph{Consequence for the proxy.} Under the median aggregator the true
influence is a step function of the local order-statistic spacing; the
proxy $\sum_i\alpha_i p_i$ is a monotone upper envelope of it, and
\eqref{eq:medianalloc} is its median-aggregator replacement.

\emph{CIFAR-10 validation.} We implement \eqref{eq:medianalloc} as an
explicit order-statistic allocator (OSA) and compare it against OURS
(Theorem~\ref{thm:t2}) on the real CIFAR-10 pipeline over
$\sigma_{\mathrm{rel}}{\in}\{0.5,1,2\}\times\kappa{\in}\{0.5,2\}$.
In the single-round scalar setting OSA attains the same median shift
$D$ as OURS to four decimals at every cell ($0.0218$ for
$\sigma_{\mathrm{rel}}\le1$, $0.0433$ at $\sigma_{\mathrm{rel}}{=}2$)
while consuming $450$--$1600\times$ less budget on CIFAR-10
($4.4$--$1600\times$ across the two real pipelines), and its damage is
\emph{budget-saturated}: granting it OURS's full budget changes $D$ by
less than $10^{-4}$.

\emph{End-to-end check.} In the full backdoor pipeline ($5$ rounds
$\times$ $5$ seeds; Table~\ref{tab:osae2e}, Fig.~\ref{fig:pareto}) OSA
is the stealthiest allocator (round detection $40$--$47\%$, $\le0.3\%$
of the nominal budget) but its end-to-end backdoor gain is statistically
null ($\Delta\mathrm{ASR}=0.004\pm0.007$ at $\kappa{=}0.5$,
$0.009\pm0.010$ at $\kappa{=}2$): it moves the median projection, not
the decision boundary, which is governed by the raw injected mass---the
client-blind W25 owns the ASR frontier of Table~\ref{tab:r2bd}. Each
allocator is thus optimal for the damage model it optimizes; no
allocator dominates across damage measures (see Discussion).

\begin{table}[!t]
\centering
\caption{Unified end-to-end comparison on the CIFAR-10 backdoor probe
(fixed budget $\Gamma=\kappa K\|\Delta\|^{2}$; means over
$\sigma_{\mathrm{rel}}{\in}\{0.5,1,2\}$, $5$ seeds per cell).
$\Delta\mathrm{ASR}$: gain over the no-attack baseline ($0.0866$);
rdet\%: round-level detection rate; $B/\Gamma$: deployed perturbation
budget as a fraction of the nominal budget (range over $\sigma$).
W25/Q25/OURS rows are the Table~\ref{tab:r2bd} means; OSA rows are a
$5$-seed rerun of the identical protocol (\S\ref{sec:median-alloc}),
whose OURS cross-check agrees within seed noise. Main-task accuracy
stays within $7.9$ pp of the no-attack baseline in the most aggressive
cell (W25 at $\sigma_{\mathrm{rel}}{=}2$, $\kappa{=}0.5$), within
$2.5$ pp for OURS, and within $0.1$ pp for OSA.}
\label{tab:osae2e}
\footnotesize\setlength{\tabcolsep}{4pt}
\begin{tabular}{lcccccc}
\toprule
& \multicolumn{3}{c}{$\kappa{=}0.5$} & \multicolumn{3}{c}{$\kappa{=}2$} \\
\cmidrule(r){2-4}\cmidrule(l){5-7}
allocator & $\Delta$ASR & rdet\% & $B/\Gamma$ & $\Delta$ASR & rdet\% & $B/\Gamma$ \\
\midrule
W25  & 0.168 & 43 & 17--100\% & 0.130 & 55 & 4--80\% \\
Q25  & 0.131 & 71 & 100\%     & 0.078 & 97 & 100\% \\
OURS & 0.062 & 75 & 100\%     & 0.040 & 84 & 47--48\% \\
OSA  & $0.004{\pm}0.007$ & 40 & $\le0.3\%$ & $0.009{\pm}0.010$ & 47 & $\le0.3\%$ \\
\bottomrule
\end{tabular}
\end{table}

%======================================================================
\section{Safety Threshold and Stackelberg Equilibrium}
\label{sec:stackelberg}
%======================================================================
\subsection{T3: phase transition}
\begin{theorem}[Safety threshold]\label{thm:t3}
With capacity \eqref{eq:etamax}, the defender-tolerable privacy level
(noise scale)
given attacker resources $(K,\Gamma)$ is
$\sigma_{\max}(K,\Gamma)=\Theta\!\bigl(\sqrt{\Gamma/K}\, d^{-1/4}\bigr)$,
\emph{decreasing} in $K$ and $d^{1/4}$: stronger privacy (larger
$\sigma$) enlarges the stealthy-attack region
$\{(K,\Gamma): K\,\eta_{\max}^2(\sigma)\ge \Gamma\}$.
\end{theorem}
\emph{Proof.} See the Supplementary Material. \hfill$\square$

Empirically, at $\rho=30$, $d=400$: separation crosses
$\mathrm{sep}^{*}$ at $\sigma\approx2.65$, consistent with the law
$\eta_{\max}(d{=}400,\sigma)\!=\!30$ (a self-consistency check, not an
independent measurement; Fig.~\ref{fig:2}).

\emph{Remark (heterogeneity floor on real pipelines).}
Theorem~\ref{thm:t3} uses the isotropic law
$\eta_{\max}\propto\sigma\, d^{1/4}$. On real pipelines the two-regime
refinement (Proposition~\ref{prop:tworegime}) replaces it by
$\eta_{\max}=\eta_{\mathrm{f}}\sqrt{1+(\sigma/\sigma_c)^2}$: in the
floor regime ($\sigma\ll\sigma_c$) the deployability condition
$K\,\eta_{\mathrm{f}}^2\ge\Gamma$ is $\sigma$-independent, so no
privacy level deploys the attack ($C<\Gamma$; the effective
$\sigma_{\max}$ is $0$ at that budget), and the $d^{1/4}$ scaling and
the $\sigma_{\max}(K,\Gamma)\propto\sqrt{\Gamma/K}\,d^{-1/4}$ law
recover only in the noise-dominated regime $\sigma\gtrsim\sigma_c$.
The practical guidance is therefore two-phase: compare $C$ against
$\Gamma$ first; only in the capacity-available phase does the privacy
level $\sigma_{\max}$ become the operative knob.

\begin{figure}[!t]
\centering
\includegraphics[width=0.6\linewidth]{figures/fig2_safety_boundary.pdf}
\caption{Safety threshold: the stealthy-escape boundary
$\eta_{\max}(\sigma)$; larger $\sigma$ enlarges the escape region
(T3).}
\label{fig:2}
\end{figure}

\subsection{T4: Stackelberg analysis}
We adopt the Stackelberg security-game paradigm
\cite{tambe2011security}: the defender picks $\sigma$ trading privacy benefit $a\ln\sigma$
against two \emph{measured} costs: the clean-utility loss of noisy
aggregation and the breach envelope $D(\sigma)$ (the best-attacker
hidden damage of \S\ref{sec:real}):
\begin{equation}
  U(\sigma)=a\ln\sigma-b_1\sigma-b_2\sigma^2-D(\sigma).
\end{equation}
Both cost terms are calibrated on the CIFAR-10 pipeline. Probing
noise-only aggregation over the same attack window gives the utility
loss $L(\sigma)=0.00278\,\sigma+0.00037\,\sigma^2$ ($R^2=0.998$), and
the R2 best-attacker envelope linearizes in the
floor regime to $D(\sigma)=0.0202+0.0109\,\sigma$ (three envelope
points; $R^2=0.93$, max residual $0.0026$---too few points for a formal
confidence interval, so its uncertainty is quantified by slope
perturbation below). The privacy
exchange rate $a$
(accuracy-equivalent value per e-fold of privacy) is a policy
parameter, not identifiable from measurements; we therefore report the
optimal response as a curve: $\sigma^\star$ solves
$2b_2\sigma^2+(b_1+d_1)\sigma-a=0$, giving $\sigma^\star=1.9$ at the
calibration $a=0.0286$, and $\sigma^\star\in[0.4,6.5]$ for
$a\in[0.005,0.12]$, always interior and well below the
crossover $\sigma_c=13.4$ of Proposition~\ref{prop:tworegime}.
Two structural facts are invariant to $a$: $U$ is strictly
concave on $\sigma>0$ ($U''(\sigma)=-a/\sigma^{2}-2b_{2}<0$), hence has
a unique interior optimum; and at $\sigma^\star$ the marginal
cost decomposes into $72\%$ breach-envelope growth versus $28\%$
direct utility loss---over-privatization is priced not in the
defender's own accuracy but in the enlarged stealthy-attack region,
the quantitative form of the T3--T4 coupling ($U$ falls by $3.2\times$
from $\sigma^\star$ to $\sigma_{\mathrm{rel}}{=}8$;
Fig.~\ref{fig:5}A).

\emph{Sensitivity and validity.} The $72/28$ split is a calibrated
value, not a universal constant: perturbing the $D$ slope by
$\pm20\%$ moves it to $66$--$76\%$ while $\sigma^\star$ stays within
$[1.7,2.2]$. The linearized $D$ is a floor-regime approximation: every
reported $\sigma^\star(a)$ lies in $[0.4,6.5]\ll\sigma_c$, inside the
window where the linearization holds, and beyond $\sigma_c$ the
saturating capacity law makes $D$ grow sublinearly, so the quadratic
utility-loss term still dominates at large $\sigma$ and $U$ retains a
unique interior maximizer. The curve $\sigma^\star(a)$---not the single
value $\sigma^\star{=}1.9$---is therefore the primary calibrated object
(Fig.~\ref{fig:5}B).

\begin{figure}[!t]
\centering
\includegraphics[width=0.8\linewidth]{figures/fig5_stackelberg.pdf}
\caption{Stackelberg analysis (T4), calibrated on CIFAR-10.
\textbf{A}: measured components of $U(\sigma)$ (privacy benefit
$a\ln\sigma$, measured utility loss, measured breach envelope
$D(\sigma)$) with interior optimum $\sigma^\star=1.9$ at the
calibration $a=0.0286$; over-privatizing beyond $\sigma^\star$
subsidizes the adversary. \textbf{B}: the optimal response
$\sigma^\star(a)$---the primary calibrated object, since $a$ is a
policy parameter---is interior and far below the two-regime crossover
$\sigma_c$ for all plausible privacy valuations; $\sigma^\star{=}1.9$
is its value at $a=0.0286$, one point on this curve.}
\label{fig:5}
\end{figure}

%======================================================================
\section{Experimental Evaluation}
\label{sec:exp}
%======================================================================
We report results at two fidelity levels. A \emph{proxy} level---the
detection-theoretic framework with large $\sigma$ and $\Gamma$ sweeps over
six allocators and mechanism-faithful baselines---validates the capacity
law and the water-filling optimizer and isolates each baseline's failure
mode; because its qualitative conclusions are not artifacts of the impact
proxy (an ablation removing capping costs $66.4\%$ on average and removing
client selection $53.1\%$; scale robustness from $(N,K){=}(20,6)$ to
$(100,20)$ leaves every conclusion unchanged), we report it in full in the
Supplementary Material and foreground the real pipeline here. The
real-pipeline evaluation asks: (i) whether the two-regime refinement of T1
survives on a real FL pipeline (R1, \S\ref{sec:r1group}); (ii) how the
stealth-constrained allocation trades off end-to-end against published
baselines, and where the aggregate-capacity condition $C/\Gamma$ places
each dataset (R2, R2b, \S\ref{sec:real}); and (iii) whether the
capacity law survives a third dataset (R5, \S\ref{sec:r5}). Robustness to
a stronger attack direction (R2c), a different aggregation rule (R3), and
an online adaptive attacker (R4) is documented in the Supplementary
Material. The labels R1--R5 name the five experimental claims and are used
for cross-reference throughout.

%----------------------------------------------------------------------
\subsection{Real pipeline I: setup and two-regime capacity
(R1)}\label{sec:r1group}
\subsubsection{Setup}
FashionMNIST with a $d{=}20490$ TinyCNN; $N{=}20$ heterogeneous
clients ($300$--$900$ samples), $K{=}6$ compromised; client-level
Gaussian DP; coordinate-median aggregation. Each attack round retrains
locally, injects along the persistent-damage direction $-\hat g$
(negative validation-gradient direction), adds DP noise, and the
server removes flagged clients (FPR $5\%$). The main comparison
(Table~\ref{tab:r2}) fixes $\sigma_{\mathrm{rel}}{=}1$ and budget
$\Gamma=\kappa\cdot K\cdot\|\Delta\|^2$ with $\kappa\in\{0.5,2\}$;
each cell averages $5$ attack rounds $\times$ $10$ seeds; per-client
capacities $\eta_{\max,i}$ are measured by bisection on the
deployment detector. Full protocol and hyperparameters are listed in the
Supplementary Material, together with the per-round
$\varepsilon(\sigma)$ accounting.


\subsubsection{R1: two-regime capacity (T1 refinement)}
As predicted by Proposition~\ref{prop:tworegime}, the measured
per-client capacity (Table~\ref{tab:r1}; Fig.~\ref{fig:7}) saturates into a heterogeneity floor
($\eta_{\max}\approx 1.85$--$2.36$ for $\sigma_{\mathrm{rel}}\le 8$)
and returns to the proportional law only beyond the fitted crossover,
with the ratio
$\eta_{\max}/\sigma$ falling from $3.69$ to a plateau at $0.15$ for
$\sigma_{\mathrm{rel}}\ge 24$---the real-pipeline two-regime law of
\S\ref{sec:t1-refine}.
\begin{figure}[!t]
\centering
\includegraphics[width=\linewidth]{figures/fig7_r1_tworegime.pdf}
\caption{R1: real-pipeline two-regime stealthy capacity (T1 refinement).
\emph{Left:} $\eta_{\max}$ vs.\ relative DP noise $\sigma_{\mathrm{rel}}$ on
both datasets, with the naive $\eta_{\max}\propto\sigma$ reference
(dashed). Low-$\sigma$ slope $\approx0.13$ (log-log), far below the
proportional reference: the capacity saturates into a heterogeneity
floor. \emph{Right:} $\eta_{\max}/\sigma$ falls with $\sigma$ and
plateaus at $\approx0.15$ for $\sigma_{\mathrm{rel}}\ge 24$: the
$\propto\sigma$ law is recovered in the noise-dominated regime.}
\label{fig:7}
\end{figure}

\begin{table}[!t]
\centering
\caption{Real-pipeline stealthy capacity $\eta_{\max}$ (in units of the
benign update norm) vs.\ relative DP noise $\sigma_{\mathrm{rel}}$ (R1).
The ratio $\eta_{\max}/\sigma$ falls from $3.69$ to $0.15$ and plateaus
there in the noise-dominated regime ($\sigma_{\mathrm{rel}}\ge 24$).}
\label{tab:r1}
\begin{tabular}{lrrrrrrrrrr}
\toprule
$\sigma_{\mathrm{rel}}$ & 0.5 & 1 & 2 & 4 & 8 & 16 & 24 & 32 & 48 & 64\\
\midrule
$\eta_{\max}$ & 1.85 & 1.93 & 2.05 & 2.20 & 2.36 & 3.02 & 3.83 & 4.97 & 7.11 & 9.75\\
$\eta_{\max}/\sigma$ & 3.69 & 1.93 & 1.02 & 0.55 & 0.29 & 0.19 & 0.16 & 0.16 & 0.15 & 0.15\\
\bottomrule
\end{tabular}
\end{table}

\subsection{Real pipeline II: allocation trade-off and the
capacity-bound regime (R2, R2b)}\label{sec:real}
\subsubsection{R2: end-to-end allocation trade-off}
On FashionMNIST ($\sigma_{\mathrm{rel}}{=}1$, baseline accuracy
$0.868$ after the same attack window) every allocator sits on the
noise floor: hidden accuracy drop $\le0.6\%$ at $\kappa{=}2$, and no
allocator's drop differs significantly from another across the $12$
$(\sigma,\kappa)$ cells (full paired-comparison statistics in the
Table~\ref{tab:r2} caption).
\emph{Two robust conclusions survive the noise floor:}
(i)~the two-regime capacity law of R1 reproduces
($\eta_{\max}/\sigma$ range $19.5$, Table~\ref{tab:r1}); and
(ii)~the aggregate stealth capacity exceeds the nominal budget
($C/\Gamma{=}8.8$, deployable; Table~\ref{tab:r2cifar}), so the
stealth-constrained water-filling optimum can deploy its full budget
on this dataset. Allocator superiority in absolute damage is not
established on FashionMNIST, however: that scale, and the allocator
frontier it makes measurable, only becomes meaningful where the task
margin is tight, which is the CIFAR-10 capacity-bound
regime of \S\ref{sec:r2b}.

\begin{table}[!t]
\centering
\caption{Real-pipeline trade-off on FashionMNIST (R2, $\sigma_{\mathrm{rel}}{=}1$,
$5$ rounds $\times$ $10$ seeds). Every allocator's hidden accuracy drop is
on the noise floor ($\le 0.6\%$ at $\kappa{=}2$); inter-allocator
differences are small (max point estimate $0.39$ pp at $\sigma{=}2$, $\kappa{=}2$)
and no allocator's hidden drop differs significantly from another across the
$12$ $(\sigma,\kappa)$ cells ($180$ paired comparisons, $\alpha{=}0.05$; median
paired $p{=}0.56$, worst $95\%$ CI covers $0$). The meaningful
damage scale and allocator frontier appear on CIFAR-10 (\S\ref{sec:r2b}).}
\label{tab:r2}
\begin{tabular}{llrr}
\toprule
strategy & $\kappa$ & acc.\ drop & mal.\ flag\%\\
\midrule
NONE (no attack) & -- & 0 & --\\
OURS & 0.5 & 0.0027 & 14.3\\
S21  & 0.5 & 0.0023 & 13.0\\
W25  & 0.5 & 0.0014 & 10.7\\
H25  & 0.5 & 0.0024 & 14.7\\
Q25  & 0.5 & 0.0019 & 13.3\\
UNI  & 0.5 & 0.0027 & 15.0\\
\midrule
OURS & 2.0 & 0.0054 & 17.0\\
S21  & 2.0 & 0.0041 & 18.7\\
W25  & 2.0 & 0.0050 & 14.0\\
H25  & 2.0 & 0.0062 & 16.0\\
Q25  & 2.0 & 0.0057 & 14.3\\
UNI  & 2.0 & 0.0052 & 16.3\\
\bottomrule
\end{tabular}
\end{table}

\subsubsection{R2b: cross-dataset capacity-bound regime (CIFAR-10)}
\label{sec:r2b}
To test whether the FMNIST frontier result transfers, we repeat R2 on
CIFAR-10 with a deeper CNN ($d{=}60362$, benign update norm
$\|\Delta\|{=}1.214$) over a $3{\times}4$ grid of
$\sigma_{\mathrm{rel}}{\in}\{0.5,1,2\}$ and
$\kappa{\in}\{0.25,0.5,1,2\}$, $5$ attack rounds $\times$ $10$ seeds
($72$ cells; Fig.~\ref{fig:8}). This grid revises the single-dataset
picture in two respects.

\emph{(i) Over-privatization helps the attacker --- confirmed on CIFAR-10 (T3/C2).}
On CIFAR-10 the upper envelope of achievable damage across all six
allocators monotonically increases with $\sigma_{\mathrm{rel}}$:
$0.0240\,{\to}\,0.0337\,{\to}\,0.0411$ at
$\sigma_{\mathrm{rel}}{=}0.5,1,2$. Larger DP noise enlarges the
masking-escape zone, so a defender who raises
$\sigma$ to ``improve privacy'' simultaneously enlarges the attacker's
stealthy-attack region. On FashionMNIST this same envelope
($0.0072\,{\to}\,0.0062\,{\to}\,0.0084$ at
$\sigma_{\mathrm{rel}}{=}0.5,1,2$) is non-monotonic and sits on the noise
floor (Table~\ref{tab:r2}); T3's \emph{magnitude} effect is therefore not
measurable on that easy task, but the capacity-regime mechanism (the
$C/\Gamma$ flip in (ii) below) carries the cross-dataset conclusion.

\emph{(ii) A capacity-bound regime where stealth-aware optimality breaks.}
The capacity regime flips between datasets. On FashionMNIST the aggregate
stealth capacity $C{=}\sum_i\eta_{\max,i}^2$ exceeds the nominal budget
($C/\Gamma{=}8.8$, deployable; Table~\ref{tab:r2cifar}) so water-filling
can deploy its full stealth-constrained budget; but because
FashionMNIST is a high-accuracy task the realized damage of every allocator
sits on the noise floor (Table~\ref{tab:r2}) and allocator superiority is
not established there. The regime where the stealth-aware optimum actually
\emph{binds} in absolute damage is the opposite, capacity-bound regime: on
CIFAR-10 the per-client capacity is anticorrelated with client
leverage: the high-data clients that carry the largest $\alpha_i$ also
have the largest benign updates and therefore the \emph{lowest}
$\eta_{\max,i}$ ($0.06$--$0.13$ vs.\ $1.6$ for small clients), giving
$C/\Gamma{=}0.48\,{<}\,1$ (capacity-bound; Fig.~\ref{fig:8}b). Water-filling then
caps every allocated client at its small $\eta_i$ and cannot deploy
the surplus budget; the entire Pareto frontier of the $72$ CIFAR runs is
owned by the uncapped/loosely-capped W25 and Q25 allocators (Fig.~\ref{fig:8}a),
and OURS is strictly interior.

Theorem~\ref{thm:t2} anticipates this outcome: the
water-filling optimum is optimal \emph{among
stealth-constrained} attackers, but when $C{<}\Gamma$ the stealth
constraint itself throttles achievable damage below what an unconstrained
attacker extracts at the price of detection. The defender's prescription
is unchanged: over-provisioning $\sigma$ is doubly harmful. It
enlarges $C$ toward $\Gamma$ (pushing the system from capacity-bound
toward deployable, where the stealth-aware attacker regains the frontier)
and enlarges the per-round escape zone.

\paragraph{$C/\Gamma$ is a threat-model quantity, not a dataset label}
Because the nominal budget is $\Gamma=\kappa K\|\Delta\|^{2}$, the ratio
$C/\Gamma$ varies as $1/\kappa$: on CIFAR-10 it reads $3.7$--$4.1$ at
$\kappa{=}0.25$ (deployable) versus $0.47$--$0.52$ at $\kappa{=}2$
(capacity-bound; Table~\ref{tab:r2cifar}), so cross-dataset comparisons
must fix $\kappa$ first, and the summation set of $C$ is the $K$
\emph{compromised} clients selected by the attacker.

Because acc-drop leaves the allocator ordering statistically unresolved in the
binding regime, we introduce a backdoor attack-success-rate probe that resolves
it in \S\ref{sec:r2bd}.

\begin{table}[!t]
\centering
\caption{CIFAR-10 cross-dataset check (R2b). Upper envelope of damage
across all allocators grows with $\sigma$ (T3/C2). The deployability
ratio $C/\Gamma$ is computed with $C{=}\sum_{i\in\mathcal{K}}\eta_{\max,i}^2$
summed over the $K{=}6$ compromised clients and
$\Gamma{=}\kappa K\|\Delta\|^2$; because $\Gamma\propto\kappa$, $C/\Gamma\propto 1/\kappa$,
so the deployable/capacity-bound label is a statement about the
\emph{assumed} attacker budget at $\kappa{=}2$, not an intrinsic
dataset property. Values: FMNIST is deployable ($C{>}\Gamma$) while
CIFAR-10 is capacity-bound ($C{<}\Gamma$) at $\kappa{=}2$.}
\label{tab:r2cifar}
\footnotesize\setlength{\tabcolsep}{3.5pt}
\begin{tabular}{lrrrr}
\toprule
dataset & $\|\Delta\|$ & $C{=}\sum_{i\in\mathcal{K}}\eta_i^2$ & $\Gamma$ ($\kappa{=}2$) & $C/\Gamma$\\
\midrule
FashionMNIST & 0.36 & 13.6 & 1.56 & 8.8\,dep.\\
CIFAR-10 & 1.214 & 8.5 & 17.7 & 0.48\,cap.\\
\bottomrule
\end{tabular}
\end{table}

\begin{figure}[!t]
\centering
\includegraphics[width=0.8\linewidth]{figures/fig8_cifar_grid.pdf}
\caption{CIFAR-10 R2 ($3{\times}4$ $\sigma$-$\kappa$ cells, $72$
runs). \textbf{a}: Stealthy-damage trade-off: the entire Pareto
frontier is owned by W25/Q25; OURS is strictly interior
(capacity-bound regime). \textbf{b}: The deployability flip:
$C/\Gamma{>}1$ on FMNIST vs.\ $C/\Gamma{<}1$ on CIFAR-10.}
\label{fig:8}
\end{figure}

\subsubsection{R2b-ASR: backdoor ASR as a finer damage probe under the stealth budget (CIFAR-10)}
\label{sec:r2bd}
The capacity-bound conclusion above is drawn from \emph{accuracy drop} as the
damage measure. Under the stealth budget acc-drop is weak: every allocator's
realized damage sits on the noise floor (Table~\ref{tab:r2},
Fig.~\ref{fig:8}a), so the ordering of allocators in the regime where
stealth-aware optimality actually binds ($\kappa{=}2$) is statistically
unresolved by acc-drop. We therefore add a finer probe whose damage is
achievable at $1$--$2\%$ of the per-update norm while leaving main-task
accuracy essentially untouched, so that ``stealth'' and ``damage'' decouple and
the allocator ordering becomes resolvable.

\emph{Setup.} Target class $0$, a fixed $3{\times}3$ pixel trigger in the
bottom-right corner, and a $1000$-image backdoor test set. The injection
direction is $v_{\mathrm{bd}}{=}-\nabla_{\theta}\mathcal{L}_{\mathrm{CE}}$
evaluated on a batch of trigger-patterned images at the target class, so the
attacker pushes compromised updates toward labeling the trigger as the target;
magnitude is capped by the same per-client $\eta_{\max,i}$ probe and the same
$\Gamma{=}\kappa K\|\Delta\|^{2}$ budget as R2/R2b. Detection uses the identical
median-score design separation $\mathrm{sep}^{*}{=}2$ ($\approx 5\%$ FPR,
$\approx 36\%$ TPR at the critical level; the realized per-allocator operating
points are in \S\ref{sec:real}) and flag-and-remove rule. The reported quantity
is the backdoor attack-success rate $\mathrm{ASR}$ and its gain over the
no-attack baseline $\Delta\mathrm{ASR}{=}\mathrm{ASR}{-}\mathrm{ASR}_{\mathrm{none}}$,
over the $\sigma_{\mathrm{rel}}{\in}\{0.5,1,2\}\times\kappa{\in}\{0.5,2\}$
grid $\times$ $5$ seeds.

\emph{Hypothesis (Theorem~\ref{thm:t2} under a sensitive measure).} Because
backdoor damage is reachable at a fraction of the norm an acc-drop needs, the
$C{<}\Gamma$ regime should resolve the allocator ordering that acc-drop
left flat: the surplus-budget allocators should lead on $\Delta\mathrm{ASR}$.
Theorem~\ref{thm:t2} predicts the uncapped allocator (Q25) and the
noise-adaptive global cap (W25) lead, while the mean-leverage water-filling
allocator (OURS), optimal only for the mean-type accuracy-drop proxy,
should fall back once the damage measure becomes order-statistic; this
is the empirical bridge to the order-statistic allocator of
\S\ref{sec:median-alloc}.

\emph{Results.} Table~\ref{tab:r2bd} reports mean $\Delta\mathrm{ASR}$ and
round detection over $\sigma_{\mathrm{rel}}{\in}\{0.5,1,2\}$ with per-cell
$95\%$ confidence intervals ($5$ seeds per cell, $df{=}4$). The probe behaves
as designed: $\Delta\mathrm{ASR}{>}0$ for every allocator at every
cell (backdoor success is achievable at every cell; the mean main-task
accuracy drop of the most aggressive cell is $7.9$ pp, $2.5$ pp for
OURS, and $0.1$ pp for OSA, Table~\ref{tab:osae2e}); stealth and damage thus decouple and the
ordering becomes resolvable. $\Delta\mathrm{ASR}$ also
rises monotonically with $\sigma_{\mathrm{rel}}$, as a larger noise budget should
permit. The finer probe resolves the ordering that acc-drop left flat, and
in a way that exposes the limit of the mean-leverage objective: the global-cap
allocator \textbf{W25} attains the largest $\Delta\mathrm{ASR}$ on average in both
$\kappa$ regimes (mean $0.168$ at $\kappa{=}0.5$, $0.130$ at $\kappa{=}2$) and
the lowest round detection ($43\%$ and $55\%$), so it Pareto-dominates
OURS on both damage and stealth. The seed-level evidence supports the
W25-over-OURS ordering specifically: the paired per-cell difference
$W25{-}OURS$ is positive in all six cells (mean $+0.04$ to $+0.17$;
CI excluding zero in the three highest-damage cells), whereas W25's
edge over Q25 is not cell-wise resolved (two of six cells reverse) and
is claimed only at the mean level. OURS itself falls to the \emph{bottom} of the
$\Delta\mathrm{ASR}$ ranking in both regimes ($0.062$ and $0.040$) while
incurring \emph{higher} detection than W25---the stealth-cap operating
point, not stealth negligence: $\eta_{\max}$ is defined at design
separation $\mathrm{sep}^{*}=2$, itself a $\approx36\%$-per-round flag
point, and an allocator that saturates caps (OURS deploys $100\%$ of
$\Gamma$ at $\kappa{=}0.5$) inherits that operating point, whereas
W25's client-blind magnitude $\min\{\kappa\sigma,\sqrt{\Gamma/K}\}$
sits below the cap for most clients. This matches
Theorem~\ref{thm:t2}'s $C{<}\Gamma$ regime (the surplus-budget allocators lead),
but it also shows that the mean-leverage water-filling objective is optimal only
for the \emph{mean-type} (accuracy-drop) damage proxy, not for the
order-statistic (backdoor/median) damage model where a simple global cap is
strictly better. That gap is what \S\ref{sec:median-alloc} closes for
the median-shift damage model; the end-to-end backdoor comparison that
includes the order-statistic allocator is reported in
Table~\ref{tab:osae2e}.

\begin{table}[!t]
\centering
\caption{CIFAR-10 backdoor ASR under the stealth budget (R2b-ASR,
$\sigma_{\mathrm{rel}}{\in}\{0.5,1,2\}$, $\kappa{\in}\{0.5,2\}$, $5$ seeds).
$\Delta\mathrm{ASR}$ is the gain over the no-attack baseline; rdet\% is the
round-level detection rate. $\pm$: largest per-cell $95\%$ CI half-width
over the three $\sigma$ cells ($df{=}4$); entries are means over
$\sigma_{\mathrm{rel}}{\in}\{0.5,1,2\}$; the full $(\sigma,\kappa)$ grid is in the
supplementary CSV.}
\label{tab:r2bd}
\footnotesize\setlength{\tabcolsep}{3.5pt}
\begin{tabular}{lrrrr}
\toprule
allocator & \multicolumn{2}{c}{$\kappa{=}0.5$} & \multicolumn{2}{c}{$\kappa{=}2$} \\
\cmidrule(r){2-3}\cmidrule(r){4-5}
 & $\Delta\mathrm{ASR}$ & rdet\% & $\Delta\mathrm{ASR}$ & rdet\% \\
\midrule
OURS & $0.062{\pm}0.084$ & 75 & $0.040{\pm}0.127$ & 84 \\
S21  & $0.087{\pm}0.053$ & 73 & $0.057{\pm}0.059$ & 81 \\
W25  & $0.168{\pm}0.073$ & 43 & $0.130{\pm}0.189$ & 55 \\
H25  & $0.088{\pm}0.027$ & 59 & $0.068{\pm}0.089$ & 83 \\
Q25  & $0.131{\pm}0.100$ & 71 & $0.078{\pm}0.088$ & 97 \\
UNI  & $0.102{\pm}0.056$ & 57 & $0.045{\pm}0.081$ & 83 \\
\bottomrule
\end{tabular}
\end{table}

Robustness to attack direction (R2c), aggregation rule (R3), and online
adaptive attackers (R4) is documented in the Supplementary Material.

\begin{figure}[!t]
\centering
\includegraphics[width=0.95\linewidth]{figures/fig13_pareto_bd.pdf}
\caption{Fixed-budget stealth--damage plane of the backdoor probe
(CIFAR-10; means over $\sigma_{\mathrm{rel}}{\in}\{0.5,1,2\}$, $5$
seeds; light points: per-$\sigma$ cells). Left: $\kappa{=}0.5$; right:
$\kappa{=}2$. W25 owns the damage frontier at both budget regimes; OSA
sits in the low-detection corner at $\Delta\mathrm{ASR}\approx0$; OURS
is dominated on both axes. Allocator optimality is damage-model
dependent (Table~\ref{tab:osae2e}).}
\label{fig:pareto}
\end{figure}

\subsection{R5: cross-dataset replication of the capacity law on MNIST}
\label{sec:r5}

The main pipeline uses FashionMNIST and CIFAR-10 (a second, color
$32{\times}32$ modality). We now ask whether the two-regime capacity law
(T1) survives a third,
independent real dataset that shares FashionMNIST's grayscale
$28{\times}28$ modality but has a distinct distribution: MNIST
($d{=}20490$, a small CNN), trained to $97.7\%$ clean accuracy versus
$61\%$ on CIFAR-10 and $84.6\%$ on FashionMNIST. The allocation-level
conclusions (T2--T3) are not re-testable on MNIST: as shown below, every
allocator sits on its noise floor.

\emph{T1 replicates.} The real-pipeline $\eta_{\max}(\sigma)$ on MNIST is
$1.43,1.53,1.62$ at $\sigma_{\mathrm{rel}}{=}0.5,1,2$; the
$\eta_{\max}/\sigma$ ratio falls from $2.86$ to $0.81$ (range $3.54$),
i.e.\ far from constant, so the heterogeneity-floor refinement (not
the naive $\eta\propto\sigma$ law) holds on MNIST exactly as on CIFAR-10
(the cross-dataset figure in the Supplementary Material, panel~A).

\emph{Absolute attack impact is task-dependent; the mechanism is not.}
Under the same R2 grid the accuracy drop of \emph{every} allocator on
MNIST is $\approx 0$ (maximum $8{\times}10^{-4}$, within seed noise),
versus $0.2$--$0.8\%$ on CIFAR-10 (Supplementary Material,
cross-dataset figure, panel~B). This is not a
theory failure: the stealth-aware optimum is a statement about the
\emph{feasibility boundary and allocation structure} (the capacity ceiling
$C/\Gamma$), not a guarantee of absolute damage. The realized damage is
absorbed by the task's own robustness: a $97.7\%$ classifier leaves
almost no room for a stealthy perturbation to move decisions, whereas a
$61\%$ one does. The structural ordering replicates: on both datasets
OURS stays within the baseline stealth envelope, and no uncapped baseline
is the unique optimum on MNIST. With R3 this corroborates, on a third
dataset and a second modality, that the capacity law and the $C/\Gamma$
boundary are robust to dataset and aggregator; the allocation-level
claims (T2--T3) rest on FashionMNIST and CIFAR-10 only.

%======================================================================
\subsection{Computational and detection overhead}
%======================================================================
A deployability diagnostic is only practical if its computational cost is
negligible relative to the FL training loop. We therefore report the
attack-side allocation solve time and the defender-side per-round cost.

\paragraph{Allocation solve time}
Table~\ref{tab:overhead} reports the median single-solve time of the four
allocators at the main client count $N{=}20$ (with $K{=}6$
compromised); the closed-form allocators stay sub-microsecond even at
$N{=}500$. Water-filling (OURS)
is a bisection over the KKT condition and is the heaviest solver, yet it
remains sub-millisecond at the deployed scale; W25/Q25 are constant-time
global scalings and OSA is a single sort. The solver is therefore not a
bottleneck---the per-round TinyCNN forward/backward pass dominates.

\begin{table}[!t]
\centering
\caption{Attack-side allocation solve time (median over $10^3$ runs, pure
Python). Water-filling (OURS) is a bisection; W25/Q25 are closed-form
global scalings; OSA is a closed-form order-statistic sort.}
\label{tab:overhead}
\footnotesize\setlength{\tabcolsep}{4pt}
\begin{tabular}{lrrrr}
\toprule
Allocator & Method & $N{=}20$ ($\mu$s) & $N{=}100$ ($\mu$s) & $N{=}500$ ($\mu$s) \\
\midrule
OURS & water-filling (bisection) & 223 & 983 & 4797 \\
W25  & global scaling (closed)   & 0.8 & 3.4 & 16 \\
Q25  & global scaling (closed)   & 1.6 & 6.3 & 31 \\
OSA  & order-statistic (sort)    & 2.9 & 13  & 70 \\
\bottomrule
\end{tabular}
\end{table}

\paragraph{Defender-side per-round cost}
The defender's detection runs a coordinate-wise median over the $N{=}20$
client updates each round---sorting $20$ scalars $d$ times, i.e.\
$O(N\,d\log N)$ with $d$ the model dimension (thousands for TinyCNN)---and
DP injection adds $N\,d$ Gaussian samples. Both are sub-millisecond per
round, so detection and aggregation, not allocation, dominate the
defender's compute; the $C/\Gamma$ check itself is a one-time offline
computation from historical updates (Implementation and Deployment
Considerations).

%======================================================================
\section{Discussion and Limitations}
%======================================================================
This section separates what the diagnostic establishes from what it
does not claim.
\begin{itemize}
  \item \textbf{A diagnostic, not a better allocator.} The contribution
  of this paper is the mechanism---the allocation optimum and the
  $C/\Gamma$ regime map---not a dominance claim for any single
  allocator: no allocator dominates across damage measures, so the
  empirical standing of OURS below is an expected consequence of the
  regime structure, not a defeated method. The stealth-aware
  optimum is optimal \emph{among stealth-constrained} attackers
  (Theorem~\ref{thm:t2}), but on neither primary dataset is its damage
  advantage over unconstrained baselines measurable. In the
  \emph{deployable} regime (FashionMNIST, and plausibly MNIST;
$C/\Gamma{>}1$) the
  task margin is wide and every allocator lands on the noise floor; in
  the \emph{capacity-bound} regime (CIFAR-10; $C/\Gamma{<}1$) the
  stealth constraint itself throttles damage, and unconstrained
  allocators (W25/Q25) own the entire Pareto frontier and
  Pareto-dominate OURS on the backdoor probe (Table~\ref{tab:r2bd}).
  The order-statistic allocator (OSA) is maximally stealthy yet its
  end-to-end backdoor gain is null (Table~\ref{tab:osae2e}); a
  deployable, margin-tight setting would be needed to observe OURS's
  frontier advantage.
  \item \textbf{Scope of validity.}\label{sec:aggr-swap} The proxy
  framework (Supplementary Material) \emph{instantiates} the capacity law
  and optimizer under the isotropic regime of T1 (a self-consistency
  check; T1 itself is established by the data collapse of
  Fig.~\ref{fig:4}); baselines are mechanism-faithful proxies, not
  official implementations, so directional conclusions are expected to
  transfer while absolute numbers will change. Theory and measurement
  center on distance-to-aggregate rules: the Supplementary Material verifies that
  Trimmed-Mean and Krum preserve the floor regime and the $C/\Gamma$
  boundary, while FLAME (cosine statistics) leaves the attack nearly
  unfiltered---a scope statement, not a rebuttal. Probing is modeled
  as a one-shot offline read; the Supplementary Material shows an online reactive
attacker gains nothing (signed gain $-0.003$ to $-0.020$). Full online
re-estimation of $\{\eta_{\max,i}\}$ remains open. All results concern the per-update DP noise model
  (\S\ref{sec:setting}).
  \item \textbf{Calibration and scope.} The Stackelberg cost terms
  ($b_1,b_2$ and the breach envelope $D$) are calibrated on the
  CIFAR-10 pipeline (\S\ref{sec:stackelberg}); the privacy exchange
  rate $a$ is a policy parameter, so we report the optimum as the
  curve $\sigma^\star(a)$ rather than a single value, and the
  $72/28$ marginal decomposition and the interior-optimum conclusion
  are the calibrated claims. We do not propose new defenses;
  defender-side mitigation is left to future work.
\end{itemize}

%======================================================================
\section{Conclusion}
%======================================================================
Attacker-side optimization reveals a quantitative cost of
over-privatization in FL: the stealthy capacity grows with $\sigma$, the
stealth-constrained allocation is water-filling (optimal for the
attacker's mean-type influence objective) with leveraged client
selection, and its deployability is bounded by the aggregate stealth
capacity $C{=}\sum_i\eta_{\max,i}^2$ relative to the nominal budget
$\Gamma$. When $C{<}\Gamma$ the stealth-aware optimum is capacity-bound
and overtaken by unconstrained allocators; the defender's utility is
non-monotone in the privacy budget in either regime. All of these
conclusions are statements about the per-update (client-level) DP noise
model, in which the detector observes noised updates; under central DP
applied only to the aggregate the masking channel vanishes
(\S\ref{sec:setting}); attacker placement is the strongest-information assumption (a
malicious or compromised server), so client-side adversaries are
conservatively upper-bounded. Two stress tests
delimit these conclusions. Under a backdoor ASR probe the allocator
ordering resolves in the binding regime and the unconstrained allocator
Pareto-dominates the water-filling optimum: the stealth-aware optimum is
optimal \emph{for its damage model}, so defenders should audit candidate
defenses under several damage measures rather than one. Under the median
aggregator, realized damage is governed by order statistics rather than by
the influence-weighted sum: a closed-form order-statistic allocation
attains the same median shift in the single-round scalar setting for
$4.4$--$1600\times$ less budget across
the two real pipelines, and its binding constraint is the benign order
statistics rather than the budget itself; at that budget scale, however,
its end-to-end backdoor gain is null, so the order-statistic objective
is the right allocator for the median-shift damage model, not for raw
attack success. Defense design should
therefore co-optimize privacy and robustness at the interior optimum
rather than max out $\sigma$. Concretely, the $C/\Gamma$ regime map
supplies a \emph{phase-conditional} prescription rather than a single
knob. In the \emph{deployable} phase ($C/\Gamma{>}1$ at the planned
$\kappa$), the binding lever is the attack surface, not noise: shrink
$K$ (client screening, attestation), tighten clipping so the per-client
caps $\eta_{\max,i}$ fall, or raise the defense budget $\kappa$---raising
$\sigma$ alone only enlarges the stealthy region while buying detection
loss. In the \emph{capacity-bound} phase ($C/\Gamma{<}1$), the stealth
constraint already throttles the attacker: tune $\sigma$ for the
privacy--utility trade-off at the interior optimum $\sigma^\star$, and
re-run the check when the pipeline, dimension, or client population
changes. Near the boundary ($C/\Gamma{\approx}1$) the phase is fragile:
re-evaluate $C/\Gamma$ for every planned $(\sigma,\kappa)$ pair,
because small changes in either flip the regime and, with it, the
defender's optimal action.

One caveat remains: the water-filling optimum and the Stackelberg
equilibrium are stated for the attacker's linear influence \emph{proxy}; under the deployed
median aggregator realized damage is governed by order statistics, so they
bound---rather than equal---the realized optimum. What we establish is
the mechanism and the $C/\Gamma$ diagnostic; we do not claim universal
optimality.

%======================================================================
\section*{Data availability}
%======================================================================
The attack--defense pipeline code, configuration files, per-run result
tables (CSV), and the automated data-consistency checker that verifies
every number reported in the figures and tables of this paper are provided
with the submission as a public repository (link in the cover
letter) and will remain publicly available upon
acceptance (repository URL/DOI to be inserted at publication).

\clearpage
\section*{Acknowledgements}
%======================================================================
This work was supported by the China University Industry--Academia
Research Innovation Fund---New Generation Information Technology
Innovation Project (Grant No.~2021ITA01022). The numerical
calculations in this paper have been done on the computing facilities
in the High Performance Computing Platform of Anqing Normal University.

%======================================================================
\section*{CRediT authorship contribution statement}
%======================================================================
\textbf{Xiaolin Zhang:} Conceptualization, Methodology, Software,
Formal analysis, Investigation, Data curation, Writing -- original draft,
Visualization. \textbf{Lei Lu:} Conceptualization, Validation, Writing --
review \& editing, Supervision, Project administration, Funding acquisition.

%======================================================================
\section*{Declaration of competing interest}
%======================================================================
The authors declare that they have no known competing financial interests
or personal relationships that could have appeared to influence the work
reported in this paper.

%======================================================================
\section*{Declaration of generative AI and AI-assisted technologies in the manuscript preparation process}
%======================================================================
During the preparation of this work the authors used AI-assisted tools in
order to assist with code development, numerical verification of reported
results, and language editing. After using these tools, the authors
reviewed and edited the content and take full responsibility for the
content of the published article.

%======================================================================
\begin{thebibliography}{37}
\bibitem{zhou2024roppfl}
H. Zhou, Y. Zheng, and X. Jia, ``Towards robust and privacy-preserving
federated learning in edge computing,'' \emph{Computer Networks},
vol.~243, art.~110321, 2024, doi:10.1016/j.comnet.2024.110321.
\bibitem{liu2025antidotefl}
Y. Liu, H. Zhang, M. Wang, Q. Xie, and Z. Sun, ``AntidoteFL: Enhancing
defense against poisoning attacks in federated learning,'' \emph{Computer
Networks}, vol.~269, art.~111427, 2025, doi:10.1016/j.comnet.2025.111427.
\bibitem{zhao2025fedmp}
K. Zhao, L. Wang, F. Yu, B. Zeng, and Z. Pang, ``FedMP: A multi-pronged
defense algorithm against Byzantine poisoning attacks in federated
learning,'' \emph{Computer Networks}, vol.~257, art.~110990, 2025,
doi:10.1016/j.comnet.2024.110990.
\bibitem{mcmahan2017}
H.~B. McMahan \etal, ``Communication-efficient learning of deep networks
from decentralized data,'' in \emph{Proc. AISTATS}, 2017.
\bibitem{kairouz2021advances}
P. Kairouz \etal, ``Advances and open problems in federated learning,''
\emph{Foundations and Trends in Machine Learning}, vol.~14, no.~1--2,
pp. 1--210, 2021.
\bibitem{li2020federated}
T. Li, A.~K. Sahu, A. Talwalkar, and V. Smith, ``Federated learning:
Challenges, methods, and future directions,'' \emph{IEEE Signal
Processing Magazine}, vol.~37, no.~3, pp. 50--60, 2020.
\bibitem{dwork2006calibrating}
C. Dwork, F. McSherry, K. Nissim, and A. Smith, ``Calibrating noise to
sensitivity in private data analysis,'' in \emph{Proc. TCC}, 2006.
\bibitem{abadi2016}
M. Abadi \etal, ``Deep learning with differential privacy,'' in
\emph{Proc. CCS}, 2016.
\bibitem{mcmahan2018dp}
H.~B. McMahan, D. Ramage, K. Talwar, and L. Zhang, ``Learning
differentially private recurrent language models,'' in \emph{Proc.
ICLR}, 2018.
\bibitem{geyer2017dp}
R.~C. Geyer, T. Klein, and M. Nabi, ``Differentially private
multi-client federated learning,'' arXiv:1712.07557, 2017.
\bibitem{blanchard2017}
P. Blanchard \etal, ``Machine learning with adversaries: Byzantine
tolerant gradient descent,'' in \emph{Proc. NeurIPS}, 2017.
\bibitem{yin2018}
D. Yin \etal, ``Byzantine-robust distributed learning: Towards optimal
statistical rates,'' in \emph{Proc. ICML}, 2018.
\bibitem{cao2021fltrust}
X. Cao, M. Fang, J. Liu, and N.~Z. Gong, ``FLTrust: Byzantine-robust
federated learning via trust bootstrapping,'' in \emph{Proc. NDSS},
2021.
\bibitem{nguyen2022flame}
T.~D. Nguyen \etal, ``FLAME: Taming backdoors in federated learning,''
in \emph{Proc. NeurIPS}, 2022.
\bibitem{mothukuri2021survey}
V. Mothukuri, R.~M. Parre\~no, S. Gharehbaghi, H.~M. Huang, and Y. Si,
``A survey on security and privacy of federated learning,'' \emph{Future
Generation Computer Systems}, vol.~115, pp. 619--640, 2021.
\bibitem{hossain2025mpelm}
M.~T. Hossain, S. Badsha, H. La, S. Islam, and I. Khalil, ``Exploiting
Gaussian noise variance for dynamic differential poisoning in federated
learning,'' \emph{IEEE Trans.\ Artificial Intelligence}, vol.~6, no.~11,
pp. 2922--2939, Nov. 2025, doi:10.1109/TAI.2025.3540030.
\bibitem{wang2025ldpfl}
Z. Wang, W. Tong, T. Han, H. Chen, T. Zhang, Y. Mao, and S. Zhong, ``On
evaluating the poisoning robustness of federated learning under local
differential privacy,'' arXiv:2509.05265, 2025. The code repository
announced in the paper (\url{https://github.com/ZiJW/LDPFL-Attack}) is
no longer publicly accessible as of Oct.\ 2026.
\bibitem{biggio2012poisoning}
B. Biggio, B. Nelson, and P. Laskov, ``Poisoning attacks against
support vector machines,'' in \emph{Proc. ICML}, 2012.
\bibitem{meizhu2015teaching}
S. Mei and X. Zhu, ``Using machine teaching to identify optimal
training-set attacks on machine learners,'' in \emph{Proc. AAAI}, 2015.
\bibitem{shejwalkar2021turning}
V. Shejwalkar and A. Houmansadr, ``Manipulating the Byzantine: Optimizing
model poisoning attacks and defenses for federated learning,'' in
\emph{Proc. NDSS}, 2021.
\bibitem{bhagoji2019analyzing}
A.~N. Bhagoji, S. Chakraborty, P. Mittal, and S. Calo, ``Analyzing
federated learning through an adversarial lens,'' in \emph{Proc. ICML},
2019.
\bibitem{fang2020local}
M. Fang, X. Cao, J. Jia, and N.~Z. Gong, ``Local model poisoning attacks
to Byzantine-robust federated learning,'' in \emph{Proc. USENIX Security
Symposium}, 2020.
\bibitem{biggio2018wild}
B. Biggio and F. Roli, ``Wild patterns: Ten years after the rise of
adversarial machine learning,'' \emph{Pattern Recognition}, vol.~84,
pp. 317--331, 2018.
\bibitem{sun2019can}
Z. Sun, P. Kairouz, and A.~T. Suresh, ``Can you really backdoor
federated learning?'' arXiv:1911.07963, 2019.
\bibitem{bagdasaryan2020backdoor}
E. Bagdasaryan, A. Veit, Y. Hua, D. Estrin, and V. Shmatikov, ``How to
backdoor federated learning,'' in \emph{Proc. AISTATS}, 2020.
\bibitem{xie2020dba}
C. Xie, K. Huang, P.-Y. Chen, and B. Li, ``DBA: Distributed backdoor
attacks against federated learning,'' in \emph{Proc. ICLR}, 2020.
\bibitem{mhamdi2018hidden}
E.~M. El Mhamdi, R. Guerraoui, and S. Rouault, ``The hidden
vulnerability of distributed learning in Byzantium,'' in
\emph{Proc. ICML}, 2018.
\bibitem{bernstein2018signsgd}
J. Bernstein, Y.-X. Wang, K. Azizzadenesheli, and A. Anandkumar,
``signSGD: With majority vote is fast and error tolerant,'' in
\emph{Proc. ICLR}, 2018.
\bibitem{xu2024minvar}
H. Xu and T. Shu, ``Defending against model poisoning attack in
federated learning: A variance-minimization approach,'' \emph{Journal
of Information Security and Applications}, vol.~82, art.~103744, 2024.
\bibitem{xu2024agnostic}
H. Xu and T. Shu, ``Attack-model-agnostic defense against model
poisonings in distributed learning,'' \emph{Journal of Information
Security and Applications}, vol.~82, 2024.
\bibitem{luong2024fedlsae}
T.~D. Luong, V.~M. Tien, N.~H. Quyen, D.~T.~T. Hien, P.~T. Duy, and
V.-H. Pham, ``Fed-LSAE: Thwarting poisoning attacks against federated
cyber threat detection system via autoencoder-based latent space
inspection,'' \emph{Journal of Information Security and Applications},
vol.~87, art.~103916, 2024.
\bibitem{bonawitz2017practical}
K. Bonawitz \etal, ``Practical secure aggregation for privacy-preserving
machine learning,'' in \emph{Proc. CCS}, 2017.
\bibitem{song2019privacy}
L. Song, R. Shokri, and P. Mittal, ``Privacy risks of securing machine
learning models against adversarial examples,'' in \emph{Proc. IEEE
Symposium on Security and Privacy}, 2019.
\bibitem{qiu2025budget}
J. Qiu, Y. Zhao, and J. Li, ``Optimal perturbation budget allocation for
data poisoning in offline reinforcement learning,'' arXiv:2512.08485,
Dec. 2025.
\bibitem{dwork2014algorithmic}
C. Dwork and A. Roth, ``The algorithmic foundations of differential
privacy,'' \emph{Foundations and Trends in Theoretical Computer
Science}, vol.~9, no.~3--4, pp. 211--407, 2014.
\bibitem{huber1964robust}
P.~J. Huber, ``Robust estimation of a location parameter,'' \emph{Annals
of Mathematical Statistics}, vol.~35, no.~1, pp. 73--101, 1964.
\bibitem{tambe2011security}
M. Tambe, \emph{Security and Game Theory: Algorithms, Deployed Systems,
Lessons Learned}. Cambridge University Press, 2011.
\end{thebibliography}

\end{document}
